\chapter{Intégration du Machine Learning}
\label{ch:ml}

\section{Motivation}
Au-delà du reporting descriptif assuré par l'entrepôt de données et le cockpit Knowage, le projet intègre quatre familles de modèles de Machine Learning afin de répondre à des besoins analytiques plus avancés~: segmenter les usagers selon leur profil de mobilité, recommander un mode de transport personnalisé, détecter les anomalies de trafic et prédire les zones à risque d'inaccessibilité. Cette dimension prédictive constitue l'apport distinctif du projet par rapport à une solution de Business Intelligence classique, en permettant aux décideurs du CETUD d'anticiper plutôt que de simplement constater.

\section{Pourquoi un entraînement direct sur les sources}
Contrairement au cockpit Knowage, qui interroge l'entrepôt de données PostgreSQL, les modèles de Machine Learning de ce projet sont entraînés et exploités directement à partir des fichiers sources nettoyés (extraits EMD et comptages de trafic), sans passer par l'entrepôt de données. Ce choix est motivé par trois raisons~:

\begin{itemize}
\item le modèle décisionnel de l'entrepôt a été conçu pour répondre à des besoins de reporting agrégé (KPI, tableaux de bord), et certaines variables individuelles nécessaires aux modèles (par exemple l'ensemble des réponses détaillées de l'EMD utilisées en feature engineering) n'y sont pas toutes représentées à la granularité requise~;
\item l'entraînement de modèles de Machine Learning nécessite des itérations rapides sur les données (sélection de variables, comparaison d'algorithmes), un cycle plus proche du travail exploratoire en Python/Colab que des cycles de chargement par lots de l'ETL~;
\item découpler les deux chaînes (décisionnelle classique vers le DW, prédictive vers les sources nettoyées) permet de faire évoluer les modèles de Machine Learning indépendamment du calendrier de mise à jour de l'entrepôt de données, sans risque de régression sur le reporting existant.
\end{itemize}

Ce choix architectural ne remet pas en cause la cohérence de la solution~: les sources utilisées par les modèles sont les mêmes données nettoyées qui alimentent par ailleurs l'ETL Talend, garantissant que les deux chaînes restent fondées sur un référentiel de données identique.

\section{Pipeline de Machine Learning}
Les quatre modèles partagent un pipeline méthodologique commun, décliné selon les spécificités de chaque tâche.

\subsection{Prétraitement}
Les valeurs manquantes sont traitées par imputation~: médiane pour les variables numériques (robuste aux valeurs extrêmes), valeur la plus fréquente pour les variables catégorielles. Les variables catégorielles sont encodées, et les variables numériques normalisées (\texttt{StandardScaler}) lorsque l'algorithme utilisé est sensible à l'échelle des variables (K-Means notamment).

\subsection{Feature engineering}
Selon le modèle, des variables dérivées sont construites~: agrégation de plusieurs indicateurs binaires en un score composite (modèle d'inaccessibilité), construction de statistiques par mode et par segment (durée, coût), ou encore calcul de scores de distance et de densité pour la détection d'anomalies spatiales.

\subsection{Sélection des variables}
La sélection des variables s'appuie sur une combinaison d'expertise métier (variables identifiées comme pertinentes lors de l'analyse des besoins, chapitre~\ref{ch:existant}) et d'analyse statistique (importance des variables calculée a posteriori sur les modèles arborescents, permettant de valider ou d'écarter certaines variables candidates).

\subsection{Entraînement}
Chaque modèle est entraîné avec une validation croisée stratifiée à 5 plis (\texttt{StratifiedKFold}) lorsque la tâche est supervisée, garantissant une estimation robuste des performances et limitant le risque de surapprentissage.

\subsection{Évaluation}
Les modèles non supervisés sont évalués par des métriques internes (coefficient de silhouette, méthode du coude). Les modèles supervisés sont évalués par exactitude, F1-score, précision, rappel et aire sous la courbe ROC (ROC AUC), avec une attention particulière à l'équilibre précision/rappel sur les classes minoritaires.

\subsection{Sauvegarde}
Chaque modèle entraîné, ainsi que les objets de prétraitement associés (\texttt{scaler}, \texttt{encoder}) et les métadonnées (liste des variables, métriques de performance), sont sérialisés au format \texttt{.pkl} via \texttt{joblib}, afin d'être chargés tels quels par le backend FastAPI sans nouvel entraînement.

\section{Présentation des modèles}

\subsection{Segmentation des usagers (K-Means)}

\paragraph{Objectif.} Regrouper les individus enquêtés en profils de mobilité homogènes, afin de permettre des analyses et des recommandations différenciées par segment.

\paragraph{Variables.} Variables sociodémographiques et de mobilité de l'enquête individus~: âge, sexe, statut d'activité, niveau d'instruction, possession du permis, fréquence d'utilisation des transports en commun, nombre de déplacements, nombre de véhicules, revenu.

\begin{figure}[H]
\centering
\includegraphics[width=0.8\textwidth]{ml/exploration_data.png}
\caption{Exploration des données individus avant segmentation}
\label{fig:exploration}
\end{figure}

\paragraph{Algorithme.} \texttt{KMeans} (scikit-learn), avec normalisation préalable des variables (\texttt{StandardScaler}) afin d'éviter qu'une variable à forte variance (le revenu, notamment) ne domine artificiellement le calcul de distance euclidienne. Le nombre optimal de clusters a été déterminé en combinant la méthode du coude (inertie intra-cluster en fonction de $k$) et le coefficient de silhouette, conduisant à retenir $k=3$.

\begin{figure}[H]
\centering
\includegraphics[width=0.7\textwidth]{ml/choix_k.png}
\caption{Détermination du nombre optimal de clusters (méthode du coude + silhouette)}
\label{fig:choix_k}
\end{figure}

\paragraph{Résultats.} Coefficient de silhouette $\approx 0{,}240$ pour $k=3$. Trois segments ont été identifiés et caractérisés (tableau~\ref{tab:segments})~: Actifs motorisés, Étudiants mobilité douce, Travailleurs informels.

\begin{table}[H]
\centering
\caption{Segments d'usagers identifiés par K-Means}
\label{tab:segments}
\begin{tabular}{p{3.5cm}p{6cm}p{3cm}}
\toprule
\textbf{Segment} & \textbf{Profil type} & \textbf{Mode dominant} \\
\midrule
Actifs motorisés & $\sim$40 ans, actif, permis, revenu élevé & Voiture / Taxi \\
Étudiants mobilité douce & $\sim$20 ans, étudiant, budget faible & Marche / TC \\
Travailleurs informels & $\sim$30 ans, actif informel, mobilité mixte & TC / Moto \\
\bottomrule
\end{tabular}
\end{table}

\begin{figure}[H]
\centering
\includegraphics[width=0.75\textwidth]{ml/heatmap_segments.png}
\caption{Caractérisation des segments par variable (heatmap)}
\label{fig:heatmap}
\end{figure}

\paragraph{Intégration.} Le modèle est exposé via l'endpoint \texttt{GET /api/segmentation/profils} et restitué dans la page \texttt{SegmentationPage} de l'espace décideurs.

\subsection{Recommandation de mode de transport (Random Forest)}

\paragraph{Objectif.} Recommander, pour un profil d'usager donné, le ou les modes de transport les plus adaptés parmi l'ensemble des modes disponibles (transport en commun, voiture, taxi/clando, marche, moto).

\paragraph{Variables.} Profil socio-démographique (âge, sexe, statut, niveau d'instruction, permis), habitudes de mobilité (fréquence d'usage des TC, nombre de déplacements, nombre de véhicules), et revenu.

\paragraph{Algorithme.} \texttt{RandomForestClassifier} (scikit-learn), 200 arbres (\texttt{n\_estimators=200}), profondeur maximale 10 (\texttt{max\_depth=10}) pour limiter le surapprentissage, pondération des classes (\texttt{class\_weight='balanced'}) pour compenser le déséquilibre entre modes, validé par validation croisée stratifiée à 5 plis.

\paragraph{Résultats.} Exactitude $\approx 80{,}15\,\%$.

\begin{figure}[H]
\centering
\includegraphics[width=0.65\textwidth]{ml/confusion_matrix.png}
\caption{Matrice de confusion du modèle de recommandation de mode}
\label{fig:confusion}
\end{figure}

\paragraph{Intégration.} En production, le modèle retourne, via l'endpoint \texttt{POST /recommander}, non pas une seule prédiction mais le \emph{top 3} des modes les plus probables avec leurs probabilités, enrichi par des statistiques de durée et de coût (quartiles Q1/médiane/Q3) calculées par mode et par segment, ce qui permet de fournir au citoyen une recommandation nuancée plutôt qu'une réponse binaire.

\subsection{Détection d'anomalies de trafic}

\paragraph{Objectif.} Détecter, parmi les comptages de trafic, des comportements anormaux~: pics de volume inattendus, schémas horaires incohérents, sites présentant un comportement différent de leurs voisins géographiques.

\paragraph{Variables.} Volumes de comptage par site et par heure, position géographique (GPS) des sites de comptage.

\paragraph{Algorithme.} Trois algorithmes complémentaires sont appliqués en parallèle puis combinés par un vote pondéré (tableau~\ref{tab:anomalies_algo}), une observation étant déclarée anormale lorsque le score combiné dépasse le seuil de 0,6. Un clustering spatial \texttt{DBSCAN} est ensuite appliqué sur les coordonnées GPS des sites détectés comme anormaux (rayon \texttt{eps} $\approx$ 500 mètres, minimum 3 sites pour former un groupe), afin de distinguer les \emph{hotspots} de congestion structurelle des incidents ponctuels isolés.

\begin{table}[H]
\centering
\caption{Algorithmes de détection d'anomalies et pondération du consensus}
\label{tab:anomalies_algo}
\begin{tabular}{p{3cm}p{2cm}p{8cm}}
\toprule
\textbf{Algorithme} & \textbf{Poids} & \textbf{Principe} \\
\midrule
Z-Score & 0,25 & Détection univariée des valeurs extrêmes par site \\
Isolation Forest & 0,40 & Détection multivariée par isolement aléatoire (arbres) \\
Local Outlier Factor (LOF) & 0,35 & Détection par densité locale (comparaison aux voisins) \\
\bottomrule
\end{tabular}
\end{table}

\paragraph{Résultats.} Cette approche par consensus réduit significativement les faux positifs par rapport à l'utilisation d'un seul algorithme, chaque méthode ayant des forces complémentaires~: le Z-Score est rapide et interprétable mais univarié, l'Isolation Forest est robuste sur des données multivariées sans hypothèse de distribution, et le LOF détecte les anomalies de densité locale invisibles aux deux autres méthodes.

\begin{figure}[H]
\centering
\includegraphics[width=0.85\textwidth]{ml/anomalies_dashboard_v2.png}
\caption{Tableau de bord de synthèse des anomalies de trafic détectées}
\label{fig:anomalies_dashboard}
\end{figure}

\paragraph{Intégration.} Les résultats sont exposés via les endpoints \texttt{GET /api/anomalies/summary} et \texttt{GET /api/anomalies/sites}, et restitués dans la page \texttt{AnomalyDashboard} avec drill-down par site de comptage.

\subsection{Prédiction du risque d'inaccessibilité}

\paragraph{Objectif.} Prédire, pour une zone ou un ménage donné, le risque d'inaccessibilité aux transports en commun, afin d'orienter les décisions d'investissement en infrastructure.

\paragraph{Variables.} Sept indicateurs terrain (distance TC, inondation, enclavement, acces sante...). Ces sept indicateurs sont combinés en un score composite, dont le seuil de classification (ménage à risque ou non) est fixé au 60\textsuperscript{e} percentile.

\paragraph{Algorithme.} Le pipeline complet comprend l'imputation -> encodage -> normalisation -> SMOTE -> validation croisée stratifiée (4 algorithmes comparés). Le modèle \textbf{XGBoost} a été retenu comme modèle final, offrant les meilleures performances sur ce jeu de données déséquilibré.

\paragraph{Résultats.} Le tableau~\ref{tab:perf_inaccess} synthétise les performances obtenues.

\begin{table}[H]
\centering
\caption{Performances du modèle de prédiction d'inaccessibilité (XGBoost)}
\label{tab:perf_inaccess}
\begin{tabular}{ll}
\toprule
\textbf{Métrique} & \textbf{Valeur} \\
\midrule
ROC AUC & 89\,\% \\
F1-score & 72\,\% \\
\bottomrule
\end{tabular}
\end{table}

\begin{figure}[H]
\centering
\includegraphics[width=0.85\textwidth]{ml/risque_inaccessibilite_ml_v2.png}
\caption{Cartographie des zones classées par niveau de risque d'inaccessibilité}
\label{fig:risque_inaccess}
\end{figure}

L'analyse de l'importance des variables (figure~\ref{fig:importance}) confirme que les variables les plus déterminantes sont : Innodations, distance à pied jusqu'à l'arrêt TC, et durée d'accès aux soins de santé.

\begin{figure}[H]
\centering
\includegraphics[width=0.75\textwidth]{ml/importance_features.png}
\caption{Importance des variables du modèle de prédiction d'inaccessibilité}
\label{fig:importance}
\end{figure}

Une analyse complémentaire par valeurs de Shapley (SHAP) a été menée afin de renforcer l'explicabilité du modèle auprès des décideurs (figure~\ref{fig:shap}), en détaillant, pour chaque prédiction individuelle, la contribution positive ou négative de chaque variable au score de risque.

\begin{figure}[H]
\centering
\includegraphics[width=0.75\textwidth]{ml/shap_inaccessibilite.png}
\caption{Analyse d'explicabilité SHAP du modèle de prédiction d'inaccessibilité}
\label{fig:shap}
\end{figure}

\paragraph{Intégration.} Le modèle est exposé via l'endpoint \texttt{POST /predict-inaccessibility}, et restitué à la fois dans la page \texttt{ZonesRisquePage} (cartographie des zones) et dans la page \texttt{SimulateurRisque}, qui permet à un décideur de faire varier les sept indicateurs pour simuler le score de risque d'une zone hypothétique.


\section{Bilan des mod\`eles de Machine Learning}

Le tableau~\ref{tab:synthese_ml} pr\'esente une synth\`ese comparative des quatre mod\`eles int\'egr\'es dans la solution.

\begin{table}[H]
\centering
\caption{Synth\`ese des mod\`eles de Machine Learning}
\label{tab:synthese_ml}
\begin{tabular}{p{3.8cm}p{2.8cm}p{4cm}p{3.5cm}}
\toprule
\textbf{Mod\`ele} & \textbf{Type} & \textbf{Objectif} & \textbf{Performance} \\
\midrule
K-Means ($k=3$) & Non supervis\'e & Segmentation des usagers & Silhouette $\approx 0{,}240$ \\
Random Forest & Supervis\'e & Recommandation de mode & Accuracy $\approx 80{,}15\,\%$ \\
Z-Score + Isolation Forest + LOF & D\'etection d'anomalies & Anomalies de trafic & Consensus pond\'er\'e $\geq 0{,}6$ -- Accord inter-mod\`eles $90{,}6\,\%$ \\
XGBoost & Supervis\'e & Pr\'ediction d'inaccessibilit\'e & ROC AUC $= 0{,}89$ \\
\bottomrule
\end{tabular}
\end{table}

\section{Intégration dans FastAPI}
Les quatre modèles, ainsi que leurs objets de prétraitement et métadonnées, sont chargés en mémoire une seule fois au démarrage du serveur FastAPI, dans un bloc de chargement global qui interrompt volontairement le démarrage si un fichier de modèle venait à manquer (afin d'éviter de servir des prédictions incohérentes). Les endpoints orchestrent ensuite, à chaque requête, la construction des features à partir des données transmises, l'appel au modèle déjà chargé, puis le formatage de la réponse au format JSON attendu par le frontend.

\section{Exploitation des prédictions}
Du côté citoyen, les prédictions du modèle de recommandation alimentent directement le parcours de planification de trajet (\texttt{FormPage}). Du côté décideurs, les prédictions des trois autres modèles sont exploitées de façon complémentaire~: la segmentation guide la priorisation des actions de communication et d'adaptation de l'offre par profil d'usager~; la détection d'anomalies alerte sur les dysfonctionnements de trafic nécessitant une intervention opérationnelle~; et la prédiction du risque d'inaccessibilité oriente les décisions d'investissement en infrastructure, avec un outil de simulation permettant d'évaluer l'impact attendu d'un projet avant sa réalisation.

\section{Valeur ajoutée de l'intelligence artificielle}
L'intégration de ces quatre modèles transforme la solution d'un outil de reporting descriptif en un outil d'aide à la décision véritablement prédictif et prescriptif~: au lieu de seulement constater des volumes de trafic ou des taux d'accessibilité passés, le système est désormais capable d'anticiper les comportements (segmentation, recommandation) et les risques (anomalies, inaccessibilité), avec un niveau de transparence (importance des variables, explicabilité SHAP) qui permet aux décideurs du CETUD d'auditer et de faire confiance aux résultats produits.

\section{Conclusion}
Ce chapitre a présenté l'intégration du Machine Learning au sein de la solution~: la justification du choix d'un entraînement direct sur les sources, le pipeline méthodologique commun aux quatre modèles, leur présentation détaillée (segmentation, recommandation, détection d'anomalies, prédiction d'inaccessibilité), leur intégration technique dans le backend FastAPI, et la valeur ajoutée qu'ils apportent aux décideurs du CETUD. Le chapitre suivant présente le développement de la plateforme applicative qui expose ces résultats aux utilisateurs.

\clearpage
